\begin{textsample}{NEG Dim 2 – mg – Score -1.00 – mg/mg\_em\_44.txt}  \label{ex:f2_neg_099}
3.1.3.1.2 Letramento \textbf{literário} Nessa perspectiva, propõe -se realizar um trabalho pedagógico com o letramento \textbf{literário} com vistas ao letramento crítico, possibilitando ao estudante, por meio das múltiplas linguagens, a interação com variados instrumentos de acesso à informação e à comunicação, em contextos de realidades físicas ou digitais.
Para além de uma breve contextualização do panorama social e histórico, que engloba as escolas \textbf{literárias} e suas características, esse componente associa -se a leituras que encaram “ os textos/enunciados como materialidades de discursos, carregados de apreciações e valores, que buscam efeitos de sentido e ecos e ressonâncias ideológicas ” ( ROJO, 2009, p. 117 ).
Nessa proposta, o estudo da literatura busca oferecer aos estudantes a possibilidade de desenvolvimento do letramento \textbf{literário}, contribuindo para sua formação integral. 3.1.3.1.3 Importância da fruição \textbf{literária} A BNCC aponta a importância da fruição \textbf{literária}, como um marco para recuperar a leitura subjetiva no Ensino Médio, ou seja, aquela que dialoga com questões internas dos sujeitos e é passível de interpretações distintas.
Tais possibilidades se traduzem numa ampliação de sua capacidade leitora de obras de outros contextos, bem como de se reconhecerem ( mais facilmente ) como sujeitos e não como receptores passivos da leitura/experiência \textbf{literária} acadêmica.
Ou seja, a interação com a multiplicidade de culturas, de textos, de contextos, de semioses em um trabalho pedagógico sistematizado contribui para a formação de leitores capazes de manipular os diferentes instrumentos multissemióticos e culturais e de utilizá -los na interação em diferentes situações comunicativas.

% matched lemmas: literário
\end{textsample}
